\paragraph{Evaluation protocol.}

To test variance correction against alternative residual models, we
compared six composition methods using a dataset of $424$ United
States equities and exchange-traded funds from Polygon.io (Table~\ref{tab:aggregate}).
Each ticker had $2{,}767$ daily closing prices from January $2014$
through December $2024$, yielding $2{,}766$ growth rates; this
complete-case selection conditioned the analysis on survival and data
coverage. We fitted the same hidden Markov generator to each asset,
including SPY (Table~\ref{tab:hyperparams}; $r_f=0$), and regressed
the other $423$ assets on observed SPY to estimate $\alpha_i$,
$\beta_i$, $R^2_{i,\real}$, and residual variance. All generators
had jumps disabled to isolate composition from jump calibration. The
naive and hybrid methods shared a centered full-return draw, with
variance correction applied only by the hybrid method
(Algorithm~\ref{alg:hybrid}). Four alternatives generated residuals
from a Gaussian distribution, a residual-fit hidden Markov model
(labeled JumpHMM-on-residuals), a stationary block bootstrap, or a
GARCH(1,1)-$t$ model, and added them to $\alpha_i+\beta_i g_m$
without centering or rescaling. Comparator fitting and eligibility
rules are given in Online Appendix~\ref{sec:supp-composition-evaluation};
$30$ non-stationary GARCH fits left $393$ eligible assets for that
method. Using observed SPY as the common market path, we ran $100$
replications per asset with seed $1234$ and pooled training scores
across tickers. We measured factor recovery by the absolute OLS
errors $|\hat\alpha-\alpha_i|$, $|\hat\beta-\beta_i|$, and
$|\hat R^2-R^2_{i,\real}|$, and naive and hybrid variance preservation
by the composed-to-generator variance ratio. Distributional scores
were KS and AD pass rates, Wasserstein-1 distance, excess kurtosis,
and the Hill index of absolute returns (Online
Appendix~\ref{sec:supp-validation}). Because only naive and hybrid
forced the residual mean to zero, a paired centering control on the
common $393$ assets separated this constraint from model fitting
(Table~\ref{tab:centering_control}; Online
Appendix~\ref{sec:supp-centering-control}).

To test performance without refitting, we froze all $2014$--$2024$
fits and evaluated $100$ simulated paths per asset against the $249$
growth rates observed in $2025$ (Table~\ref{tab:oos_composition}).
Complete holdout histories were available for $416$ assets, including
$386$ with accepted GARCH fits. Each replication used one simulated
$249$-day no-jump SPY path shared across asset compositions. The
correction used the observed training market variance in
Eqs.~\eqref{eq:scale}--\eqref{eq:r2-rescale}, while each asset's
generator variance remained draw-specific. Consequently, variation
in simulated market variance could move composed variance away from
the generator target even at zero sample cross-covariance. Observed
$2025$ returns were used only for scoring. We repeated the KS, AD,
Wasserstein-1, and kurtosis evaluations and measured variance for all
six methods relative to observed $2025$ variance. Because goodness-of-fit tests have less power on short samples, we
also compared the same simulated paths with randomly selected
contiguous $249$-day blocks of observed training returns as a
length-matched training reference. We averaged
replications within ticker before summarizing across tickers, giving
each asset equal weight.

To estimate one-day Value-at-Risk (VaR) thresholds without relying on
a single short path, we pooled daily log returns, $\Delta t$ times
the growth rates, from $5{,}000$ separately composed $249$-day paths
per asset and method (Table~\ref{tab:var}). Each replication shared
one simulated SPY path across assets and methods and followed the
holdout centering and scaling conventions. We fixed the $95\%$ and
$99\%$ coverage thresholds before scoring observed $2025$ returns;
no holdout data entered fitting or threshold calibration. These
unconditional thresholds describe a uniformly selected day within the
$249$-day construction, retain its horizon dependence, and do not
remove the zero-sum residual constraint of naive and hybrid paths.
We reported mean exceedance rates and cross-ticker standard deviations
descriptively, without treating dependent days or assets sharing a
market as independent replications. Pool-size convergence and an
independent simulation check are reported in Online
Appendix~\ref{sec:supp-var-ensemble}.


